\begin{textsample}{NEG Dim 1 – 2023\_c5 – Score 20.00 – 2023\_c5/enade\_2023\_fono\_6ca150f84b.txt}  \label{ex:f1_neg_422}
CAPÍTULO 4 DISTRIBUIÇÃO DE CONCEITOS A finalidade deste capítulo é apresentar um panorama nacional sobre a distribuição do Conceito Enade 2023.
Para traçar esse panorama, recorre -se a tabelas e gráficos, que são acompanhados de comentários e análises que contribuem para a contextualização dos dados neles expostos.
Para visualizar o panorama do Enade 2023, em nível nacional, na primeira seção deste capítulo, os dados sobre o número de cursos participantes são organizados por grande região.
Na segunda, o mapeamento do conceito do exame é feito por categoria administrativa das IES e modalidade de oferta dos cursos, de maneira cruzada às grandes regiões.
Na última seção, a lógica de apresentação dos dados é feita por organização acadêmica e, também, por grande região de funcionamento dos cursos, no país como um todo. 4.1 PANORAMA NACIONAL DA DISTRIBUIÇÃO DOS CONCEITOS A lógica adotada na apresentação dos dados, nesta seção e nas duas seguintes, pressupõe o registro de duas informações de natureza reiterativa.
O primeiro diz respeito aos valores possíveis do Enade, que variam em níveis de 1 a 5, sendo que, à medida que esses valores aumentam, o curso é melhor avaliado, e ao fato de que aqueles cursos que tiveram somente um ou nenhum concluinte participante não obtêm o Conceito Enade, ficando sem conceito ( SC ), como ressaltado na apresentação deste relatório.
O segundo registro é que, caso haja diferença entre os cursos tabulados neste capítulo e no capítulo 1, ela corresponderá aos cursos que não foram avaliados, ou seja, àqueles, em princípio, sem estudantes concluintes que tivessem sido inscritos para a prova.
Na Tabela 4.1, são apresentadas a quantidade e a distribuição de cursos de Fonoaudiologia participantes do Enade 2023, por faixa de conceito e grande região.
A partir dos dados da Tabela 4.1, nota -se que, dos 74 cursos participantes, 23 ( 31,1 \% ) classificaram -se com conceito 2, o valor modal.
Esse foi também o conceito modal nas regiões Centro-Oeste ( 25,0 \% ), Nordeste ( 33,3 \% ) e Sudeste ( 34,6 \% ).
O conceito 4 foi o segundo mais frequente em nível nacional ( 21,6 \%, correspondendo a 16 cursos ).
O conceito 3 foi o terceiro ( 20,3 \%, correspondendo a 15 cursos ).
Houve, ainda, 11 ( 14,9 \% ) cursos que receberam o conceito 5 e oito ( 10,8 \% ) cursos que receberam o conceito 1.
Dos 74 cursos de Fonoaudiologia, 1 ( 1,4 \% ) ficou sem conceito ( SC ).
A região Norte participou com 7 cursos, o que corresponde a 9,5 \% do total nacional.
Desses, três cursos receberam o conceito 1, o que equivale a 42,9 \% do total regional, o conceito modal na região.
A dois cursos ( 28,6 \% ), atribuiu -se o conceito 3 e a dois cursos ( 28,6 \% ), o conceito 2.
Ainda nessa região, nenhum curso obteve os conceitos 4 e 5 e nenhum curso ficou sem conceito ( SC ).
A região Nordeste participou com 21 cursos, correspondentes a 28,4 \% do total nacional.
Nessa região, sete cursos ( 33,3 \% em termos regionais ) obtiveram conceito 2, o conceito modal para a região, como já comentado.
Os conceitos 4 e 5 foram atribuídos, a cinco e a quatro cursos ( 23,8 \% e 19,0 \%, respectivamente ).
Ainda nessa região, três cursos ( 14,3 \% ) foram avaliados com o conceito 1 e um curso ( 4,8 \% ) foi avaliado com o conceito 3.
Um curso ( 4,8 \% ) ficou sem conceito ( SC ).
Para os 26 cursos participantes da região Sudeste ( 35,1 \% do total nacional ), o conceito 2 foi o conceito modal, como já anteriormente assinalado, tendo sido obtido por nove cursos ( 34,6 \% ).
O conceito 3 foi atribuído a oito cursos ( 30,8 \% ) e o conceito 5, a quatro cursos ( 15,4 \% ).
Receberam os conceitos 4 e 1, respectivamente, quatro cursos ( 15,4 \% ) e um curso ( 3,8 \% ).
Nessa região, nenhum curso ficou sem conceito ( SC ).
Os 16 cursos da região Sul corresponderam a 21,6 \% do total nacional.
Houve predominância do conceito 4 ( conceito modal ), atribuído a 37,5 \% dos cursos da região, ou seja, atribuído a 6 dos 16 cursos participantes na região Sul.
O conceito 2 foi atribuído a quatro cursos ( 25,0 \% ) ; o conceito 5, a três cursos ( 18,8 \% ) e o conceito 3, a três cursos ( 18,8 \% ).
Na região Sul, nenhum curso obteve conceito 1 nem ficou sem conceito ( SC ).
Já dos 4 cursos participantes na região Centro-Oeste ( 5,4 \% do total nacional ), 1 ( 25,0 \% em termos regionais ) recebeu conceito 4.
Um curso ( 25,0 \% ) obteve o conceito 3, um curso ( 25,0 \% ), o conceito 2 e um curso ( 25,0 \% ), o conceito 1.
Nessa região, nenhum curso obteve o conceito 5 e nenhum curso ficou sem conceito ( SC ).
Em um gráfico de distribuição cumulativa, em \textbf{cada} valor de abscissa ( eixo x ), o valor apresentado no eixo das ordenadas ( eixo y ) é o valor acumulado da variável ( Conceito Enade ).
Em linhas gerais, em um gráfico desta mesma tipologia, poligonais mais à esquerda representam, no caso dos dados apresentados nesta seção, as grandes regiões com a pior distribuição, e poligonais mais à direita, as grandes regiões com os melhores conceitos.
A partir destas diretrizes, no Gráfico 4.1, é apresentada a distribuição cumulativa do Conceito Enade, segundo a grande região.
No Gráfico 4.1, observando a linha vertical pontilhada, que ajuda a visualizar o percentual de conceitos acumulados até o valor 2, ou seja, conceitos 1 e 2, pode -se notar, por exemplo, que a região Norte ( linha azul ) apresenta 71,4 \% dos cursos com conceito 2 ou abaixo deste valor.
A região Sul ( linha roxa ), por outro lado, apresenta 25,0 \% dos cursos na mesma situação : com conceito 2 ou abaixo.
A região Norte ( linha azul ) apresenta os cursos com conceitos mais baixos, ao passo que a região Sul ( linha roxa ) apresenta a distribuição com valores maiores.
Os cursos em IES situadas nas regiões Centro-Oeste e Nordeste ( linhas laranja e verde, respectivamente ) apresentam uma situação intermediária. 4.2 CONCEITOS POR CATEGORIA ADMINISTRATIVA, MODALIDADE DE OFERTA DOS CURSOS E GRANDE REGIÃO Na Tabela 4.2, é apresentada a distribuição dos cursos de Fonoaudiologia participantes do Enade 2023, por categoria administrativa e modalidade de oferta, de acordo com os conceitos por eles alcançados, segundo as grandes regiões brasileiras.
De acordo com as informações da Tabela 4.2, dos 74 cursos participantes, 21 ( 28,4 \% ) eram ministrados em instituições públicas e 53 ( 71,6 \% ), em instituições privadas.
Quanto à modalidade de oferta, 72 ( 97,3 \% ) cursos eram presenciais e 2 ( 2,7 \% ), a distância.
De acordo com as informações da mesma tabela, em termos nacionais, entre instituições públicas e privadas, observa -se que, dos 11 cursos avaliados com conceito 5, 10 foram oferecidos em IES públicas, e 1, em IES privadas.
Dos 53 cursos participantes de IES privadas, o conceito 2 foi o valor modal, atribuído a 23 ( 43,4 \% ) cursos.
Em relação aos demais cursos de instituições privadas, 15 cursos ( 28,3 \% ) foram avaliados com o conceito 3, oito cursos ( 15,1 \% ) foram avaliados com o conceito 1, cinco cursos ( 9,4 \% ) foram avaliados com o conceito 4 e um curso ( 1,9 \% ), com o conceito 5.
Nessa categoria, 1 curso ( 1,9 \% ) ficou sem conceito ( SC ).
Na rede pública, o conceito modal foi o 4, atribuído a 11 cursos, correspondentes a 52,4 \% dos 21 cursos da categoria.
Dos demais cursos participantes, 10 cursos ( 47,6 \% ) obtiveram o conceito 5 e nenhum curso recebeu os conceitos 1, 2 ou 3, ou ficou sem conceito ( SC ).
Dos 72 cursos participantes da modalidade presencial, o conceito 2 foi o valor modal, atribuído a 21 ( 29,2 \% ) cursos.
Dos demais cursos dessa modalidade de oferta, 16 cursos ( 22,2 \% ) receberam o conceito 4 ; 15 cursos ( 20,8 \% ) receberam o conceito 3 ; 11 cursos ( 15,3 \% ), o conceito 5 e oito cursos ( 11,1 \% ), o conceito 1.
Nessa modalidade de oferta, um curso ( 1,4 \% ) ficou sem conceito ( SC ).
Dos 2 cursos participantes da modalidade a distância, o conceito 2 foi o valor modal, atribuído a 2 ( 100,0 \% ) cursos.
Dos demais cursos ofertados nessa modalidade, nenhum curso recebeu os conceitos 1, 3, 4 ou 5, ou ficou sem conceito ( SC ).
Na análise por região, observa -se, que, na região Norte, nenhum curso de IES públicas recebeu conceito.
As instituições privadas participaram com 7 cursos na região Norte ( 100,0 \% do total regional ) e nenhum curso ficou sem conceito ( SC ).
O conceito modal nas instituições privadas na região foi o 1, atribuído a três cursos, o que corresponde a 42,9 \% dessa categoria na região.
Os conceitos 3 e 2 foram atribuídos a dois cursos \textbf{cada} um.
Nenhum curso ficou sem conceito ( SC ) ou recebeu os conceitos 4 ou 5.
Dos cursos oferecidos na região Norte, sete eram presenciais e nenhum ficou sem conceito ( SC ).
Dos sete cursos oferecidos, três obtiveram o conceito 1 ( modal ) ; dois, o conceito 3 e dois, o conceito 2.
Nenhum curso ficou sem conceito ( SC ) ou recebeu os demais conceitos na região Norte e na modalidade presencial.
Não houve nenhum curso participante na modalidade a distância.
Na região Nordeste, a rede privada concentra 13 dos 21 cursos participantes, o equivalente a 61,9 \% do total da região.
Entre os cursos oferecidos por IES privadas no Nordeste, sete cursos foram avaliados com conceito 2, o conceito modal. três cursos obtiveram o conceito 1 ; um curso, o conceito 4 e um curso, o conceito 3.
Nessa categoria administrativa e grande região, um curso ficou sem conceito ( SC ) e nenhum curso recebeu o conceito 5.
As instituições públicas da região Nordeste participaram com oito cursos ( 38,1 \% do total da região ).
Desses, quatro ( 50,0 \% ) obtiveram conceito 5, seguidos por quatro cursos, que receberam o conceito 4.
Nenhum curso ficou sem conceito ( SC ) ou recebeu os conceitos 1, 2 ou 3.
Dos cursos oferecidos na região Nordeste, todos ( 21 ) eram presenciais : um ficou sem conceito ( SC ), três obtiveram o conceito 1 ; sete, o conceito 2 ; um, o conceito 3 ; cinco, o conceito 4 e quatro, o conceito 5.
Nenhum curso pertence à modalidade a distância.
Na região Sudeste, a proporção de cursos da rede privada, 73,1 \%, correspondeu a 19 dos 26 cursos participantes.
Nessa combinação de categoria administrativa e grande região, o conceito modal foi o 2, atribuído a nove cursos ; e nenhum curso ficou sem conceito ( SC ) ou recebeu o conceito 5.
Os demais foram avaliados com os seguintes conceitos : 3 ( oito cursos ), 4 ( um curso ) e 1 ( um curso ).
Dos sete cursos oferecidos por instituições públicas na região Sudeste, quatro obtiveram o conceito 5 e três, o conceito 4.
Nenhum curso ficou sem conceito ( SC ) ou recebeu os conceitos 1, 2 ou 3.
Ainda na região Sudeste, grande parcela dos cursos ( 96,2 \% ) foi oferecida na modalidade presencial : 25 dos 26 cursos.
Para esses, as modas foram os conceitos 2 e 3, com oito cursos \textbf{cada}.
Os demais cursos presenciais da região foram alocados nos conceitos 5 ( quatro cursos ), 4 ( quatro cursos ) e 1 ( um curso ).
Nenhum curso presencial ficou sem conceito ( SC ).
Na região Sudeste, um curso foi oferecido na modalidade a distância, que obteve o conceito 2.
As instituições privadas concentraram 11 dos 16 cursos participantes da região Sul, 68,8 \% do total regional.
Desses, quatro cursos obtiveram conceito 2, o conceito modal.
Nessa combinação de categoria administrativa e grande região, três receberam o conceito 4 ; três receberam o conceito 3 e um, o conceito 5.
Nenhum curso ficou sem conceito ( SC ) ou obteve o conceito 1.
As instituições públicas na região Sul participaram com cinco cursos ( 31,3 \% dos cursos da região ), aos quais foram atribuídos os conceitos 4 ( três cursos ) e 5 ( dois cursos ).
Nenhum curso da região Sul oferecido por IES públicas ficou sem conceito ( SC ) ou recebeu os conceitos 1, 2 ou 3.
Na região Sul, 15 dos cursos foram oferecidos na modalidade presencial.
Para esses, a moda foi o conceito 4, com 6 cursos.
Os demais cursos presenciais da região foram alocados nos conceitos 5 ( três cursos ), 3 ( três cursos ) e 2 ( três cursos ).
Nenhum curso ficou sem conceito ( SC ) ou recebeu o conceito 1.
Na modalidade a distância, foi oferecido um curso com conceito 1.
Na região Centro-Oeste, 3 dos 4 cursos participantes eram de IES privadas ( 75,0 \% em termos regionais ).
Desses, um curso recebeu conceito 1, um recebeu o conceito 2 e um, o conceito 3.
Nenhum curso recebeu os conceitos 4 ou 5 nem ficou sem conceito ( SC ).
Quanto ao único curso oferecido por instituições públicas na região Centro-Oeste ( 25,0 \% do total regional ), ele recebeu o conceito 4.
Também na região Centro-Oeste, todos os cursos foram ofertados na modalidade presencial : 4 ( 100,0 \% ).
Para esses 1 curso ( 25,0 \% ) recebeu conceito 1.
Os demais cursos presenciais da região foram alocados nos conceitos 3 ( um curso, 25,0 \% ), 2 ( um curso, 25,0 \% ) e 4 ( um curso, 25,0 \% ).
Nenhum curso presencial da região ficou sem conceito ( SC ) ou recebeu o conceito 5.
Na modalidade a distância, não foi oferecido nenhum curso.
A título de síntese, no Gráfico 4.2, apresenta -se a distribuição do Conceito Enade, segundo a categoria administrativa da IES.
Nesse gráfico, como dito anteriormente, poligonais mais à esquerda representam Conceito Enade mais baixos, e poligonais mais à direita, Conceito Enade mais altos.
É importante observar que, para \textbf{cada} poligonal presente nesse gráfico, a porcentagem constante no eixo y refere -se ao universo de cursos oferecidos por instituições públicas ( poligonal azul ) ou privadas ( poligonal verde ), respectivamente.
Os cursos em IES públicas ( poligonal azul ) apresentam uma poligonal mais à direita do que os em IES privadas ( poligonal verde ) e, consequentemente, uma distribuição de Conceito Enade com valores maiores. 4.3 CONCEITOS POR ORGANIZAÇÃO ACADÊMICA E GRANDE REGIÃO Na Tabela 4.3, encontra -se a distribuição dos conceitos atribuídos aos cursos participantes do Enade 2023, na área de Fonoaudiologia, por organização acadêmica, segundo as grandes regiões brasileiras.
Como mostram os dados da Tabela 4.3, dos 74 cursos de Fonoaudiologia participantes, 41 eram oferecidos em universidades, 26 em centros universitários e 7 em faculdades..
Essa distribuição corresponde a, respectivamente, 55,4 \%, 35,1 \% e 9,5 \% dos cursos.
Nenhum curso foi oferecido por CEFET/IF.
De acordo com os dados apresentados na mesma tabela, dos 11 cursos avaliados com o conceito 5, todos eram vinculados a universidades.
Os cursos vinculados a esse tipo de organização acadêmica tiveram o conceito 4 como conceito modal, atribuído a 13 cursos ( 31,7 \% ).
Os demais cursos ofertados por universidades receberam os conceitos 3 ( oito cursos ), 2 ( oito cursos ) e 1 ( um curso ).
Nenhum curso ficou sem conceito ( SC ).
Dos cursos vinculados a centros universitários, o conceito modal foi o 2, atribuído a 12 cursos ( 46,2 \% ).
Um curso vinculado a esse tipo de organização acadêmica ficou sem conceito ( SC ) e os outros cursos receberam os conceitos 1 ( seis cursos ), 3 ( cinco cursos ) e 4 ( dois cursos ).
Nenhum curso vinculado a centros universitários recebeu o conceito 5.
Dos 7 cursos mantidos por faculdades, 3 ( 42,9 \% ) receberam conceito 2, o conceito modal.
O segundo conceito mais frequentemente atribuído aos cursos oferecidos por essa categoria de IES foi o conceito 3, recebido por dois cursos.
Na sequência, um curso obteve conceito 4 e um curso, o conceito 1.
Dos cursos oferecidos em faculdades, nenhum curso ficou sem conceito ( SC ) ou recebeu conceito 5.
Não houve cursos oferecidos por CEFET/IF.
Considerando -se, separadamente, as regiões brasileiras, verifica -se que, na região Norte, as universidades participaram com dois cursos avaliados, aos quais foram atribuídos os seguintes conceitos : 1 a um curso e 3 a um curso.
Nenhum curso ficou sem conceito ( SC ) nem recebeu os conceitos 2, 4 ou 5.
Os centros universitários da região Norte foram representados por cinco cursos, dos quais dois cursos obtiveram conceito 1 ; dois cursos, conceito 2 e um curso, o conceito 3.
Nenhum curso ficou sem conceito ( SC ) nem recebeu os conceitos 4 ou 5.
Nem as faculdades nem os CEFET/IF participaram com cursos na região Norte.
Na região Nordeste, as universidades participaram com 12 dos 21 cursos da área de Fonoaudiologia oferecidos.
Os conceitos modais foram 5 e 4, \textbf{cada} um atribuído a quatro cursos.
Três cursos receberam o conceito 2 e um curso, o conceito 3.
Nenhum curso ficou sem conceito ( SC ) ou recebeu o conceito 1.
Os centros universitários contaram com sete cursos participantes na região Nordeste, três dos quais receberam o conceito modal 2.
Foi atribuído o conceito 1 a dois cursos e o 4 a um curso.
Um curso oferecido por centros universitários no Nordeste ficou sem conceito ( SC ) e nenhum curso recebeu os conceitos 3 ou 5.
As faculdades foram representadas por dois cursos na região Nordeste, os quais receberam os conceitos 1 ( um curso ) e 2 ( um curso ).
Nenhum curso ficou sem conceito ( SC ) nem recebeu os conceitos 3, 4 ou 5.
Na região Nordeste, os CEFET/IF não participaram com nenhum curso.
Na região Sudeste, as universidades concentraram 15 dos 26 cursos de Fonoaudiologia da região.
Dos cursos oferecidos em IES com esse tipo de organização acadêmica, nessa região, o conceito modal foi 3, atribuído a cinco cursos.
Os demais cursos receberam os conceitos 5 ( quatro cursos ), 4 ( quatro cursos ) e 2 ( dois cursos ).
Nenhum curso ficou sem conceito ( SC ) ou recebeu o conceito 1.
Os centros universitários participaram com nove cursos na região Sudeste, dos quais seis obtiveram o conceito modal 2.
O conceito 3 foi atribuído a dois cursos e o conceito 1, a um curso.
Nenhum curso ficou sem conceito ( SC ) nem recebeu os conceitos 4 ou 5 nessa combinação de organização acadêmica e região.
Já as faculdades foram representadas por dois cursos na região Sudeste, dos quais nenhum ficou sem conceito ( SC ) nem recebeu os conceitos 1, 4 ou 5.
Os conceitos 2 e 3 foram atribuídos a 1 curso \textbf{cada} um.
Os CEFET/IF não participaram com nenhum curso na região Sudeste.
Dos 16 cursos da região Sul, 10 eram mantidos por universidades, para os quais o conceito modal foi 4, atribuído a 4 cursos.
Os demais cursos receberam os conceitos 5 ( três cursos ), 2 ( dois cursos ) e 3 ( um curso ).
Em IES com esse tipo de organização acadêmica, nenhum curso da região Sul ficou sem conceito ( SC ) ou recebeu o conceito 1.
Dos três cursos participantes de centros universitários da região Sul, nenhum ficou sem conceito ( SC ) ou recebeu os conceitos 1 ou 5.
Um curso recebeu conceito 4, um curso, o conceito 3 e um curso, o conceito 2.
Quanto aos três cursos vinculados a faculdades na região Sul, um curso recebeu o conceito 4.
Os demais conceitos ficaram assim distribuídos : 3 ( um curso ) e 2 ( um curso ).
Dos três cursos, nenhum ficou sem conceito ( SC ) nem recebeu os conceitos 1 ou 5.
Na região Centro-Oeste, 2 dos 4 cursos oferecidos eram mantidos por universidades.
Para esse tipo de organização acadêmica, o conceito 4 foi atribuído a um curso.
Em relação às universidades, nenhum curso da região ficou sem conceito ( SC ) nem recebeu os conceitos 1, 3 ou 5.
Um curso recebeu o conceito 2.
Os centros universitários da região Centro-Oeste contaram com dois cursos : um curso foi avaliado com o conceito 3 e um, com o conceito 1.
Nenhum curso ficou sem conceito ( SC ) nem recebeu os conceitos 2, 4 ou 5.
Nenhum curso foi oferecido em faculdades nem em CEFET/IF na região Centro-Oeste.
De maneira semelhante à forma como se procedeu em relação à seção anterior, esta seção encerra -se com o Gráfico 4.3, que apresenta a distribuição do Conceito Enade dos cursos de Fonoaudiologia segundo a organização acadêmica da IES.
No Gráfico 4.3, apresenta -se a distribuição do Conceito Enade segundo a organização acadêmica da IES.
A leitura desse último gráfico é semelhante à dos anteriores, seguindo -se a lógica de que poligonais deslocadas para a esquerda correspondem aos conceitos mais baixos ( 1, 2 e 3 ) e poligonais deslocadas para a direita correspondem a Conceitos Enade mais altos.
Nas universidades ( linha azul ), todos os cursos obtiveram conceito 5 ou abaixo disso, já que a poligonal cumulativa alcança 100 \% para o conceito 5.
Além disso, os cursos de universidades apresentam uma poligonal mais à direita do que as demais, denotando uma melhor distribuição de Conceitos Enade.
No outro extremo, encontram -se os cursos mantidos por faculdades ( linha vermelha ), apresentando as respectivas poligonais mais à esquerda e, consequentemente, uma distribuição de Conceitos Enade com os menores valores.
Os cursos mantidos por centro universitário ( linha verde ) apresentou poligonal mais central e, consequentemente, uma distribuição de Conceitos Enade com os valores intermediários.

% matched lemmas: cada
\end{textsample}
